\exercisesfor{2}

\begin{enumerate}
\def\labelenumi{\arabic{enumi}.}
\tightlist
\item
  Let g be a free Lie algebra with basic family \((x_{1},\ldots,x_{n},\ldots)\) . Let \(g_{n}\) denote the subalgebra of g generated by \(x_{1},\ldots,x_{n}\) and let \(b_{n}\) denote the smallest ideal of \(g_{n}\) containing \(x_{n}\) .
\end{enumerate}

\begin{enumerate}
\def\labelenumi{(\alph{enumi})}
\tightlist
\item
  Show that \(\mathfrak{b}_n\) is the submodule of \(\mathfrak{g}_n\) generated by the
\end{enumerate}

\[
\left(\operatorname{ad} \left(x _ {i _ {1}}\right) \circ \dots \circ \operatorname{ad} \left(x _ {i _ {k}}\right)\right) \left(x _ {n}\right),
\]

where \(k \geqslant 0\) and \(i_h \leqslant n\) for all \(h\) .

\begin{enumerate}
\def\labelenumi{(\alph{enumi})}
\setcounter{enumi}{1}
\item
  Show that \(\mathfrak{g}_n = \mathfrak{g}_{n-1} \oplus \mathfrak{b}_n\) ; deduce that \(\mathfrak{g}\) is the direct sum of the \(\mathfrak{b}_n\) for \(n \geqslant 1\) .
\item
  Show, using (a) and (b), that the K-module g is generated by the elements \([x_{i_{1}}, [x_{i_{2}}, \ldots, [x_{i_{k-1}}, x_{i_{k}}] \ldots]]\) , where \(k \geqslant 0\) and \(i_{h} \leqslant i_{k}\) for h \textless{} k.
\end{enumerate}

¶ 2. Let X be a countable set with at least two elements and let H be the set of subsets of M(X) which are Hall sets (cf.~Definition 2). Show that Card(H) = \(2^{N_0}\).

\begin{enumerate}
\def\labelenumi{\arabic{enumi}.}
\setcounter{enumi}{2}
\item
  Let \(\mathbf{X}\) be a set with a total ordering. Show that there exists a Hall set relative to \(\mathbf{X}\) such that \(\mathrm{H} \cap \mathrm{M}^3(\mathbf{X})\) consists of the products \(z(yx)\) with \(y < x\) , \(y \leqslant z\) and that \(\mathrm{H} \cap \mathrm{M}^4(\mathbf{X})\) consists of the products \(w(z(yx))\) with \(w \geqslant z \geqslant y\) , \(y < x\) and the products \((ab)(cd)\) with \(a < b\) , \(c < d\) and either \(a < c\) or \(a = c\) and \(b < d\) .
\item
  \begin{enumerate}
  \def\labelenumii{(\alph{enumii})}
  \tightlist
  \item
    Show that the Lie algebra defined by the presentation
  \end{enumerate}
\end{enumerate}

\[
\{x, y; [ x, [ x, y ] ] = [ y, [ x, y ] ] = 0 \}
\]

admits the basis \((x,y,[x,y])\) . Define an embedding of this algebra in a matrix algebra.

\begin{enumerate}
\def\labelenumi{(\alph{enumi})}
\setcounter{enumi}{1}
\tightlist
\item
  Do the same for the Lie algebra defined by the presentation
\end{enumerate}

\[
\{x, y; [ x, [ x, y ] ] - 2 x = [ y, [ x, y ] ] + 2 y = 0 \}.
\]

\begin{enumerate}
\def\labelenumi{(\alph{enumi})}
\setcounter{enumi}{2}
\tightlist
\item
  Show that the Lie algebra with presentation
\end{enumerate}

\[
\{x, y, z; [ x, y ] - x = [ y, z ] - y = [ z, x ] - z = 0 \}
\]

reduces to 0.

\begin{enumerate}
\def\labelenumi{\arabic{enumi}.}
\setcounter{enumi}{4}
\item
  Let \(\mathfrak{g}\) be a free Lie algebra and let \(\mathfrak{r}\) be an ideal of \(\mathfrak{g}\) . Suppose that \(\mathfrak{r} = [\mathfrak{g},\mathfrak{r}]\) . Show that \(\mathfrak{r} = \{0\}\) (use the fact that \(\bigcap_{n\geqslant 0}\mathcal{C}^n\mathfrak{g} = \{0\}\) ).
\item
  Let \(\mathbf{E}\) be a module. Let \(\mathbf{ME}\) be the free algebra of \(\mathbf{E}\) (Algebra, Chapter III, §2, Exercise 13). Recall that \(\mathrm{ME} = \bigoplus_{n\geqslant 1}\mathrm{E}_n\) , where
\end{enumerate}

\[
\mathrm{E} _ {1} = \mathrm{E}, \qquad \mathrm{E} _ {n} = \bigoplus_ {p + q = n} \mathrm{E} _ {p} \otimes \mathrm{E} _ {q} \quad \text {if} \quad n \geqslant 2,
\]

and the algebra structure on ME is defined by means of the canonical mappings \(\mathbf{E}_n\otimes \mathbf{E}_m\to \mathbf{E}_{n + m}\) .

\begin{enumerate}
\def\labelenumi{(\alph{enumi})}
\item
  Let J be the smallest two-sided ideal of ME containing the elements of the form xx and \((xy)z + (yz)x + (zx)y\) with x, y, z in E. Let LE = ME/J. Show that J is a graded ideal of ME; deduce a graduation \((\mathrm{L}^{n}\mathrm{E})_{n\geqslant1}\) of LE.
\item
  Show that LE is a Lie algebra; it is called the free Lie algebra of the module E. Show that for every Lie algebra g and every linear mapping f: \(E \to g\) there exists one and only one Lie algebra homomorphism F: \(LE \to g\) which extends f.
\item
  Define isomorphisms \(\mathbf{E} \to \mathbf{L}^1\mathbf{E}\) and \(\wedge^2\mathbf{E} \to \mathbf{L}^2\mathbf{E}\) . Show that \(\mathbf{L}^3\mathbf{E}\) is identified with the quotient of \(\mathbf{E} \otimes \wedge^2\mathbf{E}\) by the submodule generated by the elements
\end{enumerate}

\[
x \otimes (y \wedge z) + y \otimes (z \wedge x) + z \otimes (x \wedge y) \quad \text { for } x, y, z \text { in } E.
\]

\begin{enumerate}
\def\labelenumi{(\alph{enumi})}
\setcounter{enumi}{3}
\item
  Show that, when E is a free module of basis X, LE is identified with the free Lie algebra \(\mathrm{L}(\mathbf{X})\) defined in no. 2.
\item
  Show that the canonical mapping of E into the enveloping algebra U(LE) of LE can be extended to an isomorphism of the tensor algebra TE onto U(LE).
\item
  Let \(\sigma\) be the linear mapping of \(\wedge^2\mathrm{E}\) into \(\mathsf{T}^2\mathsf{E}\) such that
\end{enumerate}

\[
\sigma (x \wedge y) = x \otimes y - y \otimes x.
\]

Construct a module E for which \(\sigma\) is not injective. Deduce that for this module the canonical mapping of LE into U(LE) is not injective (compare with Exercise 9 of Chapter I, §2).

\begin{enumerate}
\def\labelenumi{\arabic{enumi}.}
\setcounter{enumi}{6}
\tightlist
\item
  Suppose that K is a field. Let l be a free Lie algebra with basic family \((x_{i})_{i\in I}\) and let M be an l-module.
\end{enumerate}

\begin{enumerate}
\def\labelenumi{(\alph{enumi})}
\item
  Show that \(\mathbf{H}^2 (\mathbf{l},\mathbf{M}) = \{0\}\) , cf.~Chapter I, § 3, Exercise 12. (Use Corollary 2 to Proposition 1 and part (i) of the Exercise in question.)
\item
  Show that, for every family \((m_{i})_{i\in I}\) of elements of M, there exists a cocycle \(\phi:l\to M\) of degree 1 such that \(\phi (x_i) = m_i\) and that this cocycle is unique. Deduce an exact sequence:
\end{enumerate}

\[
0 \rightarrow \mathrm{H} ^ {0} (\mathfrak {l}, \mathrm{M}) \rightarrow \mathrm{M} \rightarrow \mathrm{M} ^ {\mathrm{I}} \rightarrow \mathrm{H} ^ {1} (\mathfrak {l}, \mathrm{M}) \rightarrow 0.
\]

If I is finite of cardinal \(n\) and M is of finite rank over K, then

\[
\operatorname{rg} \mathrm{H} ^ {1} (\mathfrak {l}, \mathrm{M}) - \operatorname{rg} \mathrm{H} ^ {0} (\mathfrak {l}, \mathrm{M}) = (n - 1) \operatorname{rg} \mathrm{M}.
\]

\begin{enumerate}
\def\labelenumi{\arabic{enumi}.}
\setcounter{enumi}{7}
\tightlist
\item
  Suppose that K is a field. Let l be a free Lie algebra with basic family \((\overline{x}_{i})_{i\in I}\) , let r be an ideal of l contained in {[}l, l{]} and let g = l/r; let \(x_{i}\) denote the image of \(\overline{x}_{i}\) in g.
\end{enumerate}

Show the equivalence of the following properties:

\begin{enumerate}
\def\labelenumi{(\roman{enumi})}
\item
  g is a free Lie algebra.
\item
  g admits \((x_{i})\) as basic family (i.e.~r = \{0\}).
\item
  For every g-module M, H²(g, M) = \{0\}.
\item
  If g operates trivially on K, then \(H^{2}(g, K) = \{0\}\) .
\end{enumerate}

(The implications (ii) \(\Rightarrow\) (i) and (iii) \(\Rightarrow\) (iv) are obvious and (i) \(\Rightarrow\) (iii) follows from Exercise 7. To prove that (iv) \(\Rightarrow\) (ii) it suffices to show that \(\mathfrak{r} = [\mathfrak{l},\mathfrak{r}]\) , cf.~Exercise 5; otherwise take a hyperplane \(\mathfrak{h}\) of \(\mathfrak{r}\) containing \([\mathfrak{l},\mathfrak{r}]\) and note that the extension \(\mathfrak{r} / \mathfrak{h}\to \mathfrak{l} / \mathfrak{h}\to \mathfrak{l} / \mathfrak{r} = \mathfrak{g}\) is essential.)

¶ 9. Let \(\mathfrak{g} = \bigoplus_{n \geqslant 1} \mathfrak{g}_n\) be a graded Lie algebra. If \(\mathfrak{h}\) is a graded subalgebra of \(\mathfrak{g}\) such that \(\mathfrak{g} = \mathfrak{h} + [\mathfrak{g}, \mathfrak{g}]\) , show that \(\mathfrak{h} = \mathfrak{g}\) . If \((x_i)\) is a family of homogeneous elements of \(\mathfrak{g}\) , deduce that \((x_i)\) is a generating family if and only if the images of the \(x_i\) in \(\mathfrak{g}/[\mathfrak{g}, \mathfrak{g}]\) generate \(\mathfrak{g}/[\mathfrak{g}, \mathfrak{g}]\) .

If K is a field on which g operates trivially and \(\mathrm{H}^{2}(\mathfrak{g},\mathbf{K})=\{0\}\) , show that the family \((x_{i})\) is basic if and only if the images of the \(x_{i}\) in g/{[}g, g{]} form a basis of this vector space. (Use Exercise 8.)†

\begin{enumerate}
\def\labelenumi{\arabic{enumi}.}
\setcounter{enumi}{9}
\tightlist
\item
  Let \(t\) be a free Lie algebra admitting a finite basic family with \(n\) elements and let \(u\) be a surjective endomorphism of \(l\) . Show that \(u\) is an isomorphism. (Write \(\mathrm{gr}^n(l) = \mathcal{C}^n / \mathcal{C}^{n+1}l\) and denote by \(\mathrm{gr}^n(u)\) the endomorphism of \(\mathrm{gr}^n(l)\) derived from \(u\) ; note that \(\mathrm{gr}^n(u)\) is surjective; as \(\mathrm{gr}^n(l)\) is a free K-module of finite rank, deduce that \(\mathrm{gr}^n(u)\) is bijective since the kernel of \(u\) is contained in \(\bigcap_{n} \mathcal{C}^n l\) , which is \(\{0\}\) .)
\end{enumerate}

If \((y_{1},\ldots ,y_{m})\) is a generating family of l, show that \(m\geqslant n\) and that we have equality if and only if \((y_{1},\ldots ,y_{m})\) is a basic family.

¶ 11. Let Y and Z be two disjoint sets (with Y totally ordered), let \(X = Y \cup Z\) and \(l = L(X)\) . Let \(a_{Y}\) be the submodule of l generated by Z and {[}l, l{]}; it is an ideal of the Lie algebra l. We propose to find explicitly a basic family of \(a_{Y}\) .

Let \(\mathbf{M}\) be the set of positive integral functions on \(\mathbf{Y}\) of finite support. If \(m\in \mathbf{M}\) , we denote by \(\theta^m\) the endomorphism of l given by

\[
\theta^ {m} = \prod_ {y \in Y} (a d y) ^ {m (y)},
\]

where the product is relative to the order relation given on Y. Let \(\mathfrak{S}\) denote the subset of \(\mathbf{Y} \times \mathbf{M}\) consisting of the ordered pairs \((u, m)\) such that there exists \(y \in \mathbf{Y}\) for which \(y > u\) and \(m(y) \geqslant 1\) . Show that the elements

\[
\theta^ {m} (u) \quad \text { for } \quad (u, m) \in \mathfrak {S} \quad \text { and } \quad \theta^ {m} (z) \quad \text { for } \quad (z, m) \in Z \times M
\]

form a basic family of the Lie algebra \(\mathfrak{a}_{\mathbf{Y}}\) .

(Reduce the problem to the case where Y is finite and argue by induction on Card(Y). Use the Corollary to Proposition 10 applied to the greatest element y of Y and apply the induction hypothesis to \(Y = \{y\}\) .)

\begin{enumerate}
\def\labelenumi{\arabic{enumi}.}
\setcounter{enumi}{11}
\tightlist
\item
  Let \(\mathbf{X} = \{x, y\}\) be a set of two elements. Show that the derived algebra of \(\mathbf{L}(\mathbf{X})\) is a free Lie algebra admitting as basic family the elements \(((\operatorname{ad}y)^q \circ (\operatorname{ad}x)^p)(y)\) , for \(p \geqslant 1\) , \(q \geqslant 0\) . (Use Exercise 11.)
\end{enumerate}

¶ 13. (In this exercise we assume that every projective module over K is free; this is so for example if K is a principal ideal domain.)

Let \(\mathfrak{l} = \sum_{n=1}^{\infty} \mathfrak{l}_n\) be a graded Lie algebra admitting a basic family B consisting of homogeneous elements and let \(\mathfrak{h}\) be a graded Lie subalgebra of \(\mathfrak{l}\) , which is a direct factor of \(\mathfrak{l}\) considered as a module.

\begin{enumerate}
\def\labelenumi{(\alph{enumi})}
\tightlist
\item
  For all \(i \geqslant 0\) , let \(\mathfrak{l}^{(i)}\) be the graded subalgebra of \(\mathfrak{l}\) such that
\end{enumerate}

\[
\begin{array}{l l} \mathfrak {l} _ {j} ^ {(i)} = \mathfrak {h} _ {j} & \text { if } j \leqslant i. \\ \mathfrak {l} _ {j} ^ {(i)} = \mathfrak {l} _ {j} & \text { if } j > i. \end{array}
\]

Then \(l = \ell^{(0)} \supset \ell^{(1)} \supset \cdots \supset \mathfrak{h}\) . Show, arguing by induction on \(i\) , the existence of a basic family \(B^{(i)}\) of \(\ell^{(i)}\) , consisting of homogeneous elements, such that the elements of \(B^{(i-1)}\) and of \(B^{(i)}\) of degree \(< i - 1\) are the same. (Suppose that \(B^{(i-1)}\) has been constructed. Let \(m_i\) be the intersection of \(l_i = l_i^{(i-2)}\) with the subalgebra generated by the \(b_j\) for \(j < i\) and let \(b_i\) be the submodule of \(l_i\) generated by the elements of \(B^{(i-1)}\) of degree \(i\) . The induction hypothesis implies that \(l_i = m_i \oplus b_i\) . As \(m_i \subset b_i\) , we can decompose \(b_i\) as \(b_i = y_i \oplus b_i\) , so that \(b_i = m_i \oplus b_i\) ; by the hypothesis on K the modules \(y_i\) and \(b_i\) are free;

changing the elements of \(\mathrm{B}^{(i-1)}\) of degree i if necessary, we can assume that \(b_{i}\) is generated by a subset of \(\mathrm{B}^{(i-1)}\) . Applying Exercise 11 to \(\mathfrak{l}^{(i-1)}\) and its ideal \(\mathfrak{l}^{(i)}\) , we obtain a basic family \(\mathrm{B}^{(i)}\) of \(\mathfrak{l}^{(i)}\) with the desired properties.)

\begin{enumerate}
\def\labelenumi{(\alph{enumi})}
\setcounter{enumi}{1}
\tightlist
\item
  Deduce from (a) the fact that \(\mathfrak{h}\) has a basic family consisting of homogeneous elements.
\end{enumerate}

¶ 14. Suppose that K is a field. Let X be a set and x, y two elements of L(X) which are linearly independent over K. Show that the family \((x,y)\) is free in the Lie algebra L(X). (Let \(x_{p}\) (resp. \(y_{q}\) ) be the non-zero homogeneous component of x (resp. y) of highest degree. Adding if necessary a multiple of y to x, we can assume that \(x_{p}\) and \(y_{q}\) are linearly independent. The subalgebra of L(X) generated by \(x_{p}\) and \(y_{q}\) is graded and hence free, cf.~Exercise 13. Deduce that \((x_{p},y_{q})\) is a free family and pass from this to \((x,y)\) .)

¶ 15. Let \(\mathrm{X} = \{x, y\}\) be a set with two elements and let \(\sigma\) be an automorphism of the Lie algebra \(\mathrm{L}(\mathrm{X})\) . Show that, if \(\mathrm{K}\) is a field, \(\sigma\) preserves the graduation of \(\mathrm{L}(\mathrm{X})\) (apply Exercise 14 to the non-zero homogeneous components of \(\sigma(x)\) and \(\sigma(y)\) of highest degree); deduce an isomorphism of \(\operatorname{Aut}(\mathrm{L}(\mathrm{X}))\) onto \(\mathbf{GL}(2, \mathrm{K})\) . Extend these results to the case where the ring \(\mathrm{K}\) has no nilpotent element \(\neq 0\) . When \(\mathrm{K}\) contains an element \(\varepsilon \neq 0\) of zero square, show that \(x \mapsto x, y \mapsto y + \varepsilon[x, y]\) can be extended to an automorphism of \(\mathrm{L}(\mathrm{X})\) which does not preserve the graduation of \(\mathrm{L}(\mathrm{X})\) .

¶ 16. Suppose that K is a Q-algebra. If g is a Lie algebra, we denote by U(g) its enveloping algebra, by σ the canonical mapping of g into U(g) and by S(g) the symmetric algebra of the K-module g. There exists one and only one linear mapping

\[
\eta (\mathfrak {g}): \mathrm{S} (\mathfrak {g}) \rightarrow \mathrm{U} (\mathfrak {g})
\]

such that \(\eta (\mathfrak{g})(x^n) = \sigma (x)^n\) for all \(x\in \mathfrak{g}\) and all \(n\in \mathbf{N}\) . Then

\[
\eta (\mathfrak {g}) (x _ {1} \dots x _ {n}) = \frac {1}{n !} \sum_ {s \in \mathfrak {S} _ {n}} \sigma (x _ {s (1)} \dots x _ {s (n)}), \quad x _ {i} \in \mathfrak {g}.
\]

We propose to show that \(\eta(g)\) is bijective.

\begin{enumerate}
\def\labelenumi{(\alph{enumi})}
\item
  Prove that \(\eta(g)\) is surjective and that it is bijective when the K-module \(g\) is free (use the Poincaré-Birkhoff-Witt Theorem).
\item
  Let \(\mathfrak{h}\) be an ideal of \(\mathfrak{g}\) . Let \(\mathfrak{s}_{\mathfrak{h}}\) (resp. \(\mathfrak{u}_{\mathfrak{h}}\) ) denote the ideal of \(S(\mathfrak{g})\) (resp. the two-sided ideal of \(U(\mathfrak{g})\) ) generated by \(\mathfrak{h}\) (resp. \(\sigma(\mathfrak{h})\) ). The diagram
\end{enumerate}

\[
\begin{array}{c} 0 \to s _ {\mathfrak {h}} \to S (g) \to S (g / \mathfrak {h}) \to 0 \\ \downarrow^ {\eta (g)} \quad \downarrow^ {\eta (g / \mathfrak {h})} \\ 0 \to u _ {\mathfrak {h}} \to U (g) \to U (g / \mathfrak {h}) \to 0 \end{array}
\]

is commutative and its rows are exact. Deduce that the image \(\tilde{u}_{\mathfrak{h}}\) of \(s_{\mathfrak{h}}\) under \(\eta(g)\) is contained in \(u_{\mathfrak{h}}\) and that \(\tilde{u}_{\mathfrak{h}} = u_{\mathfrak{h}}\) when \(\eta(g)\) and \(\eta(g/h)\) are injective (in particular, when the module \(g\) and \(g/h\) are free).

\begin{enumerate}
\def\labelenumi{(\alph{enumi})}
\setcounter{enumi}{2}
\item
  Take \(\mathfrak{g}\) to be a free Lie algebra with basic family \((\mathbf{X}_1, \ldots, \mathbf{X}_n, \mathrm{H})\) , where \(n \geqslant 0\) , and take \(\mathfrak{h}\) to be the ideal of \(\mathfrak{g}\) generated by H. Show that \(\mathfrak{g}\) and \(\mathfrak{g} / \mathfrak{h}\) are free modules. Deduce that, in that case, \(\tilde{\mathfrak{u}}_{\mathfrak{h}} = \mathfrak{u}_{\mathfrak{h}}\) . In particular, there exists \(x \in \mathfrak{s}_{\mathfrak{h}}\) such that \((\eta(\mathfrak{g}))(x) = \sigma(\mathbf{X}_1) \ldots \sigma(\mathbf{X}_n) \sigma(\mathrm{H})\) .
\item
  We return to the general case. Show that \(\mathfrak{u}_{\mathfrak{h}}\) is generated over \(\mathrm{K}\) by the elements of the form \(\sigma(x_1) \ldots \sigma(x_n)\sigma(h)\) , with \(n \geqslant 0\) , \(x_i \in \mathfrak{g}\) and \(h \in \mathfrak{h}\) (note that \(\mathfrak{u}_{\mathfrak{h}}\) coincides with the left ideal generated by \(\mathfrak{h}\) ). Show, using (c), that such an element belongs to \(\tilde{\mathfrak{u}}_{\mathfrak{h}}\) (use a suitable homomorphism of a free Lie algebra into \(\mathfrak{g}\) ). Deduce that \(\tilde{\mathfrak{u}}_{\mathfrak{h}} = \mathfrak{u}_{\mathfrak{h}}\) .
\item
  Show that, if \(\eta(g)\) is bijective, so is \(\eta(g/h)\) . Deduce finally that, for every Lie algebra \(g, \eta(g)\) is bijective (write \(g\) as the quotient of a free Lie algebra) and that the canonical homomorphism
\end{enumerate}

\[
\omega : \mathsf {S} (\mathfrak {g}) \to \operatorname{gr} \mathrm{U} (\mathfrak {g}) \quad (\text { cf.   Chapter   I,   } \S 2, \text {   no.   } 6)
\]

is an isomorphism (``Poincaré-Birkhoff-Witt Theorem'' for Lie algebras over Q-algebras).

